\documentclass{article}
\usepackage{amsmath}
\usepackage{amssymb}
\usepackage{amsthm}
\usepackage{mathtools}
\usepackage{geometry}
\usepackage{hyperref}

\newcommand{\rk}{\operatorname{rank}}
\newcommand{\GL}{\operatorname{GL}}
\newcommand{\R}{\mathbb{R}}
\newcommand{\Mat}{\mathrm{Mat}}
\newcommand{\rad}{\operatorname{rad}}

\theoremstyle{definition}
\newtheorem{theorem}{Theorem}[section]
\newtheorem{proposition}[theorem]{Proposition}
\newtheorem{lemma}[theorem]{Lemma}
\newtheorem{corollary}[theorem]{Corollary}
\newtheorem{definition}[theorem]{Definition}
\newtheorem{remark}[theorem]{Remark}

\title{Task 3: Classification of the Indecomposables}
\date{}

\begin{document}
\maketitle

\section{Goal}

This file classifies the indecomposable modules of the weight-space algebra
\[
A_L^{\mathrm{wt}}.
\]

The expected answer is a hybrid one:

\begin{itemize}
\item the non-wrapping family is the usual type-$A$ interval family on the open
  chain
  \[
  0\to 1\to \cdots \to L;
  \]
\item the wrapped family consists of the new indecomposables that pass through
  the closing arrow and are determined by the non-uniform Kupisch series
  \[
  (L+1,L,\dots,L).
  \]
\end{itemize}

\section{Interval notation}

\begin{definition}[Oriented interval]
For vertices $i,j\in\{0,\dots,L\}$, let
\[
[i\to j]_0
\]
denote the oriented interval obtained by starting at $i$, following the
orientation of the cycle, and stopping at $j$, with no repeated vertices.

Let
\[
[i\to j]_1
\]
denote the set of arrows appearing in this oriented path.

We write
\[
\delta(i,j)
\]
for the oriented distance from $i$ to $j$ along this interval, i.e.
\[
\delta(i,j):=\# [i\to j]_1.
\]
Thus
\[
0\le \delta(i,j)\le L.
\]
\end{definition}

\begin{definition}[Interval module]
For a pair $(i,j)$ with $0\le \delta(i,j)\le L$, define the interval module
$M(i,j)$ by
\[
M(i,j)_t
=
\begin{cases}
k, & t\in [i\to j]_0,\\
0, & t\notin [i\to j]_0,
\end{cases}
\]
and for each arrow $a_s:s-1\to s$ define
\[
M(i,j)_{a_s}
=
\begin{cases}
\mathrm{id}_k, & a_s\in [i\to j]_1,\\
0, & a_s\notin [i\to j]_1.
\end{cases}
\]
\end{definition}

\begin{remark}
The support interval $[i\to j]$ may or may not pass through the closing arrow.
The wrapped/non-wrapped distinction is therefore built into the index pair
$(i,j)$.
\end{remark}

\section{The admissible index set}

\begin{definition}[Admissible interval indices]
Let
\[
\mathcal I_L^{\mathrm{wt}}
:=
\{(i,j)\mid 0\le i\le j\le L\}
\cup
\{(i,j)\mid 1\le i\le L,\ 0\le j\le i-2\}.
\]
\end{definition}

\begin{remark}[Equivalent distance formulation]
Let
\[
(c_0,\dots,c_L)=(L+1,L,\dots,L)
\]
be the Kupisch series from Task 2.
Then
\[
(i,j)\in \mathcal I_L^{\mathrm{wt}}
\quad\Longleftrightarrow\quad
\delta(i,j)\le c_i-1.
\]

Equivalently:

\begin{itemize}
\item if $i=0$, then all intervals $[0\to j]$ with $0\le j\le L$ are allowed;
\item if $1\le i\le L$, then every interval starting at $i$ is allowed except
  the forbidden full wrapped interval $[i\to i-1]$ of cyclic length $L$.
\end{itemize}
\end{remark}

\begin{proposition}[Admissibility criterion for interval modules]
The interval representation $M(i,j)$ factors through the algebra
\[
A_L^{\mathrm{wt}}=kC_{L+1}/I_L^{\mathrm{wt}}
\]
if and only if
\[
(i,j)\in \mathcal I_L^{\mathrm{wt}}.
\]
\end{proposition}

\begin{proof}
Assume first that
\[
0\le i\le j\le L.
\]
Then $[i\to j]_1$ is contained in the open chain
\[
\{a_{i+1},a_{i+2},\dots,a_j\}\subseteq \{a_1,\dots,a_L\},
\]
so the closing arrow $a_0$ acts by zero on $M(i,j)$. Every relation path
$p_\ell$ with $1\le \ell\le L$ passes through $a_0$, hence each $p_\ell$ acts
by zero. Therefore $M(i,j)$ is an $A_L^{\mathrm{wt}}$-module.

Now assume that
\[
1\le i\le L,
\qquad
0\le j\le i-2.
\]
Then $[i\to j]_1$ is the wrapped path
\[
\{a_{i+1},\dots,a_L,a_0,a_1,\dots,a_j\},
\]
whose length is
\[
\delta(i,j)=L+j-i+1\le L-1.
\]
Any nonzero path acting on $M(i,j)$ must be entirely supported on the arrow set
$[i\to j]_1$. Since every relation path $p_\ell$ has length $L$, no $p_\ell$
can be contained in $[i\to j]_1$, so all of them act by zero. Hence
$M(i,j)$ again factors through $A_L^{\mathrm{wt}}$.

Finally, let
\[
1\le i\le L
\]
and consider the forbidden full wrapped interval
\[
[i\to i-1].
\]
Its arrow set is exactly the relation path $p_i$, so on the one-dimensional
vertex space at $i$ the composition
\[
p_i:M(i,i-1)_i\to M(i,i-1)_{i-1}
\]
is the identity. Thus $p_i$ acts nontrivially, and $M(i,i-1)$ does not define
an $A_L^{\mathrm{wt}}$-module.
\end{proof}

\section{The two families}

\begin{proposition}[The non-wrapping family]
If
\[
0\le i\le j\le L,
\]
then $M(i,j)$ is exactly the usual indecomposable interval module of the
equioriented type-$A_{L+1}$ quiver
\[
0\to 1\to \cdots \to L.
\]
\end{proposition}

\begin{remark}
This is the part of the classification that should be described using the same
language as Simon's type-$A$ decomposition.
\end{remark}

\begin{proposition}[Projective-quotient realization of admissible intervals]
For every admissible pair
\[
(i,j)\in \mathcal I_L^{\mathrm{wt}},
\]
there is an isomorphism
\[
M(i,j)\cong P(i)/\rad^{\delta(i,j)+1}P(i).
\]
\end{proposition}

\begin{proof}
By Task 2, the indecomposable projective $P(i)$ is uniserial of length
\[
c_i=
\begin{cases}
L+1, & i=0,\\
L, & 1\le i\le L.
\end{cases}
\]
Thus for every admissible pair, equivalently for every pair satisfying
\[
\delta(i,j)\le c_i-1,
\]
there is a unique indecomposable quotient of $P(i)$ of length
\[
\delta(i,j)+1.
\]

Under the quiver description, this quotient keeps precisely the first
$\delta(i,j)+1$ simples encountered along the oriented cycle starting at $i$.
Therefore it has one-dimensional vertex spaces exactly on the vertices of
$[i\to j]_0$, and the arrows in $[i\to j]_1$ act as identities while all other
arrows act by zero. This is exactly the interval module $M(i,j)$.
\end{proof}

\begin{remark}[The wrapped family]
If
\[
1\le i\le L,
\qquad
0\le j\le i-2,
\]
then the interval $[i\to j]$ passes through the closing arrow. These are the
genuinely cyclic indecomposables. By the proposition above, they are precisely
the nontrivial quotients of the projectives $P(i)$, $i\ge 1$, that wrap once
through the arrow $a_0$ but stop before the forbidden full wrapped interval.
\end{remark}

\section{Classification theorem}

\begin{theorem}[Indecomposables of the weight-space algebra]
Every finite-dimensional indecomposable $A_L^{\mathrm{wt}}$-module is
isomorphic to a unique interval module
\[
M(i,j)
\qquad
((i,j)\in \mathcal I_L^{\mathrm{wt}}).
\]
\end{theorem}

\begin{proof}[Proof route]
Since $A_L^{\mathrm{wt}}$ is Nakayama, every indecomposable module is uniserial
and is a quotient of a unique indecomposable projective $P(i)$.

If the top simple is at vertex $i$, then the indecomposable must be a quotient
of $P(i)$, and its composition length can be any integer between $1$ and the
projective length $c_i$. By the previous proposition, these quotients are
exactly the interval modules
\[
M(i,j)
\qquad
((i,j)\in \mathcal I_L^{\mathrm{wt}}).
\]

More concretely:
\begin{itemize}
\item for $i=0$, the admissible quotients are
  \[
  M(0,0),M(0,1),\dots,M(0,L);
  \]
\item for $1\le i\le L$, the admissible quotients are the non-wrapping modules
  \[
  M(i,i),M(i,i+1),\dots,M(i,L),
  \]
  together with the wrapped modules
  \[
  M(i,0),M(i,1),\dots,M(i,i-2).
  \]
\end{itemize}

The only excluded endpoint is $j=i-1$, which would require the forbidden full
wrapped interval.

Uniqueness follows because the top simple determines the start vertex $i$, and
the composition length determines the unique endpoint $j$ along the oriented
interval starting at $i$.
\end{proof}

\begin{corollary}[Explicit decomposition of the indecomposable family]
The indecomposables split as
\[
\{M(i,j)\mid 0\le i\le j\le L\}
\sqcup
\{M(i,j)\mid 1\le i\le L,\ 0\le j\le i-2\}.
\]
\end{corollary}

\section{Deliverables}

All Task 3 deliverables have been achieved in this file:

\begin{itemize}
\item achieved: the exact admissible index set
  $\mathcal I_L^{\mathrm{wt}}$ is defined;
\item achieved: the type-$A$ non-wrapping family is isolated explicitly;
\item achieved: the wrapped family is identified and described as projective
  quotients;
\item achieved: every indecomposable is proved to be a unique admissible interval
  module $M(i,j)$.
\end{itemize}

Therefore Task 3 is complete and the indecomposable family needed for the
Krull--Schmidt decomposition is now fixed.

\end{document}
